% =========================================================
% macros.tex  (layout inspired by the provided sample pages)
% Target compiler: XeLaTeX
% =========================================================

% -------------------------
% 0) Extra packages not in preamble_1
% -------------------------
% preamble_1 already loads: fontspec, xcolor, tikz, tcolorbox, tabularx,
%   enumitem, fancyhdr, titlesec, ifthen, amsmath, hyperref, pgfplots
% Only load what preamble_1 does NOT provide:
\RequirePackage{graphicx}
\RequirePackage{calc}
\RequirePackage{eso-pic}
\RequirePackage{etoolbox}
\usetikzlibrary{positioning,shapes.geometric,shapes.misc,shapes.symbols}
\tcbuselibrary{skins,breakable,theorems}

% -------------------------
% 1) Global look & feel
% -------------------------
% Color palette (close to the screenshots)
\definecolor{GDBlue}{HTML}{127BBE}      % cyan/blue for headings
\definecolor{GDDeepBlue}{HTML}{0D67B5}  % strong blue bars
\definecolor{GDPink}{HTML}{C1007A}      % magenta ribbon
\definecolor{GDOrange}{HTML}{C8741D}    % orange exercise bar
\definecolor{GDPurple}{HTML}{6D2C91}    % purple "Dạng" box
\definecolor{GDTeal}{HTML}{0E8C8C}      % teal section number
\definecolor{GDLiteBlue}{HTML}{EAF5FF}  % light blue box background
\definecolor{GDLiteGray}{HTML}{F5F7FA}

% Fonts and spacing are defined in preamble_1 — format_2 does not override them.
% List spacing inherits from preamble_1's parskip + setstretch.
\setlist[itemize]{leftmargin=*, itemsep=\GDListItemSep, topsep=\GDListTopSep}
\setlist[enumerate]{leftmargin=*, itemsep=\GDListItemSep, topsep=\GDListTopSep}

% Global box sizing & spacing defaults
\tcbset{
  enhanced,
  breakable,
  % rounded corners like the sample pages
  rounded corners,
  arc=\GDBoxArc,
  boxrule=\GDDefaultBoxRule,
  colback=white,
  left=\GDBoxPadH,right=\GDBoxPadH,top=\GDBoxPadV,bottom=\GDBoxPadV,
}

% Layout preset hooks (no-ops; spacing is owned by preamble_1)
\newcommand{\GDLayoutDefault}{}
\newcommand{\GDLayoutTight}{}

% -------------------------
% 2) Header / footer (book-like)
% -------------------------
\newcommand{\GDSetHeader}[2]{%
  % #1: left header text, #2: right header text
  \pagestyle{fancy}
  \fancyhf{}

  % Top blue line with rounded corner feel
  \fancyhead[C]{%
    \begin{tikzpicture}[baseline=(n.base)]
      \node[inner sep=0pt, outer sep=0pt] (n) {\phantom{X}};
      \draw[GDDeepBlue, line width=\GDHeaderRuleWidth, rounded corners=6pt]
        (0,0.22) -- (0.65\linewidth,0.22)
        .. controls +(0.8,0) and +(-0.8,0) .. (0.95\linewidth,0.22)
        -- (\linewidth,0.22);
    \end{tikzpicture}%
  }

  \fancyhead[L]{\scriptsize\color{GDDeepBlue}#1}
  \fancyhead[R]{\scriptsize\color{GDDeepBlue}#2\;\textbf{\thepage}}

  % Bottom bar
  \fancyfoot[C]{%
    \begin{tikzpicture}[baseline=(n.base)]
      \node[inner sep=0pt, outer sep=0pt] (n) {\phantom{X}};
      \draw[GDDeepBlue, line width=\GDHeaderRuleWidth] (0,0.22) -- (\linewidth,0.22);
    \end{tikzpicture}%
  }
}

% A convenience default header (override in your main file if you want)
\GDSetHeader{Chương 1. MẪU CHUYÊN ĐỀ}{Trang}

% -------------------------
% 3) Title styles
% -------------------------
\titleformat{\section}
  {\Large\bfseries\color{GDTeal}}
  {}{0em}{}
\titleformat{\subsection}
  {\large\bfseries\color{GDDeepBlue}}
  {}{0em}{}

% -------------------------
% 4) Top chapter banner + lesson line
% -------------------------
% \ChapterBanner{1}{MẪU CHUYÊN ĐỀ SỐ 2, MẪU CHUYÊN ĐỀ SỐ 2}
\newcommand{\ChapterBanner}[2]{%
  \begin{tikzpicture}
    \fill[GDDeepBlue] (0,0) rectangle (\linewidth,0.25);
    \fill[GDDeepBlue] (0,0.25) rectangle (0.66\linewidth,0.34);

    \node[anchor=west, text=white, font=\bfseries\large] at (0.02\linewidth,0.12) {\strut};

    \node[anchor=east, text=GDDeepBlue, font=\bfseries\Large] at (\linewidth, -0.15) {\strut};

    % Right "Chương" label
    \node[anchor=east, text=GDBlue, font=\bfseries\itshape\LARGE] at (\linewidth,0.55) {Ch\'\uong\;#1};
    % Main title centered
    \node[anchor=north, text=GDBlue, font=\bfseries\Huge, align=center] at (0.70\linewidth,0.28) {#2};

    % Decorative slashes
    \fill[GDBlue] (0.57\linewidth,0.02) rectangle (0.60\linewidth,0.22);
    \fill[GDBlue!60] (0.61\linewidth,0.02) rectangle (0.63\linewidth,0.22);
  \end{tikzpicture}
  \vspace{\ChapterBannerBelowSkip}
}

% \LessonLine{Bài số}{1}{BẤT PHƯƠNG TRÌNH BẬC HAI MỘT ẨN}
\newcommand{\LessonLine}[3]{%
  \begin{tikzpicture}
    \node[anchor=west, draw=GDPink!80, fill=GDPink!8, minimum height=0.85cm,
          minimum width=\linewidth, inner xsep=10pt, inner ysep=6pt] (bar) at (0,0) {};

    % left tag
    \node[anchor=west, fill=GDOrange, text=white, font=\bfseries, rounded corners=2pt,
          inner xsep=8pt, inner ysep=3pt] at ($(bar.west)+(0.18cm,0)$) {#1};

    % red circle number
    \node[circle, fill=red!85!black, text=white, font=\bfseries, inner sep=1.8pt]
      at ($(bar.west)+(1.62cm,0)$) {#2};

    % main title
    \node[anchor=west, text=GDPurple, font=\bfseries\Large]
      at ($(bar.west)+(2.25cm,0)$) {#3};

    % right slashes
    \fill[GDPurple!55] ($(bar.east)+(-1.35cm,-0.18cm)$) rectangle ($(bar.east)+(-1.05cm,0.18cm)$);
    \fill[GDPurple!30] ($(bar.east)+(-0.95cm,-0.18cm)$) rectangle ($(bar.east)+(-0.65cm,0.18cm)$);
    \fill[GDPurple!15] ($(bar.east)+(-0.55cm,-0.18cm)$) rectangle ($(bar.east)+(-0.25cm,0.18cm)$);
  \end{tikzpicture}
  \vspace{\LessonLineBelowSkip}
}

% -------------------------
% 5) Section ribbon (A/B/C/D blocks)
% -------------------------
% \SectionRibbon{A}{Tóm tắt lí thuyết}
\newcommand{\SectionRibbon}[2]{%
  \vspace{\SectionRibbonAboveSkip}
  \noindent
  \begin{tikzpicture}
    \node[anchor=west, fill=GDPink, text=white, font=\bfseries\Large, inner xsep=10pt, inner ysep=6pt] (tag) at (0,0) {#1};
    \fill[GDPink!50] ($(tag.east)+(0.12cm,-0.22cm)$) rectangle ($(tag.east)+(0.48cm,0.22cm)$);
    \fill[GDPink!25] ($(tag.east)+(0.55cm,-0.22cm)$) rectangle ($(tag.east)+(0.92cm,0.22cm)$);
    \node[anchor=west, text=GDPink, font=\bfseries\Large] at ($(tag.east)+(1.05cm,0)$) {#2};
  \end{tikzpicture}
  \vspace{\SectionRibbonBelowSkip}
}

% -------------------------
% 6) Utility commands
% -------------------------
\newcommand{\DK}{\textbf{ĐKXĐ}}
\newcommand{\Buoc}{\textbf{Bước}}

% Auto-wrap 3 columns helper
\newcommand{\ThreeColTab}[1]{%
\small
\setlength{\tabcolsep}{\ThreeColTabSep}%
\renewcommand{\arraystretch}{1.28}%
\begin{tabularx}{\linewidth}{@{}>{\raggedright\arraybackslash}X >{\raggedright\arraybackslash}X >{\raggedright\arraybackslash}X@{}}%
#1
\end{tabularx}%
\normalsize
}

\newcommand{\m}[1]{\(\displaystyle #1\)}

% -------------------------
% 7) Tcolorbox tag helper
% -------------------------
% Small inline tag used in box titles
\newcommand{\GDTag}[2]{%
  \tikz[baseline=(n.base)]\node (n) [fill=#1, text=white, font=\bfseries\footnotesize,
    rounded corners=2pt, inner xsep=\GDTagPadX, inner ysep=\GDTagPadY] {#2};%
}

% -------------------------
% 7.1) Page background / watermark (optional)
% -------------------------
% Full-page background image (e.g., watermark mascot). Usage:
%   \GDSetWatermark{watermark.pdf}{0.10}{0.85}
% You can pass pdf/png/jpg. For emf, convert to pdf first.
\newcommand{\GDSetWatermark}[3]{%
  % #1 file, #2 opacity (0..1), #3 scale
  \AddToShipoutPictureBG*{%
    \AtPageCenter{%
      \begin{tikzpicture}
        \node[opacity=#2] {\includegraphics[scale=#3]{#1}};
      \end{tikzpicture}%
    }%
  }%
}

% Remove watermark if needed
\newcommand{\GDClearWatermark}{\ClearShipoutPictureBG}

% -------------------------
% 7.2) Small numbered subheader ("1 Phần nội dung thứ nhất")
% -------------------------
% Usage: \GDSubHeader{1}{Phần nội dung thứ nhất}
\newcommand{\GDSubHeader}[2]{%
  \vspace{\GDSubHeaderAboveSkip}\noindent
  \begin{tikzpicture}
    \node[rounded corners=2.5pt, fill=GDTeal, text=white, font=\bfseries, inner xsep=10pt, inner ysep=4pt] (n) {#1};
    \node[anchor=west, font=\bfseries\large, text=GDTeal] at ($(n.east)+(8pt,0)$) {#2};
  \end{tikzpicture}\par\vspace{\GDSubHeaderBelowSkip}
}

% -------------------------
% 8) Theorem-like boxes (Định nghĩa/Định lí/Hệ quả/Nhận xét/Ghi nhớ)
% -------------------------
% Usage: \begin{gdDefinition}{1.1} ... \end{gdDefinition}
\newtcolorbox{gdDefinition}[1]{
  colframe=red!55!black,
  colback=red!3,
  title={\GDTag{red!70!black}{\,\large$\lightning$\;}\;\textbf{Định nghĩa\;#1}},
}

\newtcolorbox{gdTheorem}[1]{
  colframe=blue!60!black,
  colback=blue!3,
  title={\GDTag{blue!70!black}{\,\large$\lightning$\;}\;\textbf{Định lí\;#1}},
}

\newtcolorbox{gdCorollary}[1]{
  colframe=green!55!black,
  colback=green!3,
  title={\GDTag{green!60!black}{\,\large$\lightning$\;}\;\textbf{Hệ quả\;#1}},
}

\newtcolorbox{gdRemark}[1]{
  colframe=ForestGreen!55!black,
  colback=ForestGreen!3,
  title={\GDTag{ForestGreen!70!black}{\,\large$\lightning$\;}\;\textbf{Nhận xét\;#1}},
}

\newtcolorbox{gdWarning}{
  colframe=red!65!black,
  colback=red!2,
  title={\GDTag{red!75!black}{\Large$\triangle$}\;\textbf{Lưu ý}},
}

\newtcolorbox{gdMemory}[1]{
  colframe=GDTeal!70!black,
  colback=GDTeal!6,
  title={\GDTag{GDTeal!80!black}{\Large$\checkmark$}\;\textbf{GHI NHỚ\;--\;#1}},
}

% -------------------------
% 9) Exercise / Example boxes ("Ví dụ", "Bài", "Câu")
% -------------------------
% Usage: \begin{gdExample}{5} ... \end{gdExample}
\newtcolorbox{gdExample}[1]{
  colframe=GDBlue!90!black,
  colback=GDLiteBlue,
  boxrule=\GDExampleBoxRule,
  title={\GDTag{GDBlue!90!black}{\,\small$\diamond$\;}\;\textbf{Ví dụ\;#1}},
}

\newtcolorbox{gdPractice}[1]{
  colframe=GDOrange!85!black,
  colback=GDOrange!4,
  boxrule=\GDExampleBoxRule,
  title={\GDTag{GDOrange!85!black}{\,\small$\diamond$\;}\;\textbf{Bài\;#1}},
}

\newtcolorbox{gdQuestion}[1]{
  colframe=GDPink!85!black,
  colback=GDPink!3,
  boxrule=\GDExampleBoxRule,
  title={\GDTag{GDPink!85!black}{\,\small$\diamond$\;}\;\textbf{Câu\;#1}},
}

% -------------------------
% 9.1) Problem bar (like "Bài 2" line in the sample)
% -------------------------
% Usage:
%   \begin{gdProblemBar}{Bài 2}
%     Nội dung bài ... \ABCD{A}{...}{B}{...}{C}{...}{D}{...}
%   \end{gdProblemBar}
\newtcolorbox{gdProblemBar}[1]{
  enhanced,
  breakable,
  colback=white,
  colframe=GDOrange!80!black,
  boxrule=\GDExampleBoxRule,
  left=\GDBoxPadH,right=\GDBoxPadH,top=\GDProblemBarPadV,bottom=\GDProblemBarPadV,
  arc=\GDBoxArc,
  title={\GDTag{GDOrange!85!black}{\,\small$\diamond$\;}\;\textbf{#1}},
}

% -------------------------
% 10) Multiple-choice bubbles + answer mark
% -------------------------
\newcommand{\GDCirc}[1]{%
  \tikz[baseline=(n.base)]\node (n) [circle, draw=GDPink, very thick, text=GDPink,
  inner sep=1.2pt, font=\bfseries\small] {#1};%
}

% \ABCD{A}{...}{B}{...}{C}{...}{D}{...}
\newcommand{\ABCD}[8]{%
\begin{tabularx}{\linewidth}{@{}X X@{}}
  \GDCirc{#1}\;#2 & \GDCirc{#3}\;#4\\[2pt]
  \GDCirc{#5}\;#6 & \GDCirc{#7}\;#8\\
\end{tabularx}
}

% Mark chosen answer: \ChonDapAn{A}
\newcommand{\ChonDapAn}[1]{\textit{Chọn đáp án }\GDCirc{#1}}

% -------------------------
% 11) "Lời giải" heading + dotted writing lines
% -------------------------
\newcommand{\LoigiaiTitle}{%
  \vspace{\LoigiaiTitleAboveSkip}
  \begin{center}
    {\color{GDBlue}\bfseries\large \textbf{\raisebox{-0.2ex}{\large$\bullet$}\;Lời giải.}}
  \end{center}
}

% \DottedLines{5}  -> 5 dotted lines for students to write
\newcommand{\DottedLines}[1]{%
  \par\noindent
  \foreach \i in {1,...,#1}{%
    \noindent\tikz{\draw[GDBlue!70, dotted, line width=\DottedLineWidth] (0,0) -- (\linewidth,0);}\\[\DottedLineGap]
  }%
}

% Convenience environment:
% \begin{gdSolution}[5] ... \end{gdSolution}
\newenvironment{gdSolution}[1][4]{%
  \def\gdSolutionLines{#1}%
  \LoigiaiTitle
}{%
  \vspace{-\GDSolutionVAdjust}
  \ifnum\gdSolutionLines>0\DottedLines{\gdSolutionLines}\fi
}

% -------------------------
% 12) "Dạng" box (purple)
% -------------------------
% \begin{gdDang}{2}{Giải phương trình mũ cơ bản, phương pháp đưa về cùng cơ số} ...
\newtcolorbox{gdDang}[2]{
  colframe=GDPurple,
  colback=white,
  boxrule=\GDDangBoxRule,
  title={\GDTag{GDPurple}{DẠNG\;#1}\;\textbf{#2}},
}

% =========================================================
% 13) Extra polish (rounded header corners, checkbox, spacing)
%    + Auto-numbered environments (easy mode)
% =========================================================

% -------------------------
% 13.1) Header corner decorations + hand icon
% -------------------------
% A tiny "hand" icon (unicode). Works best with XeLaTeX.
\newcommand{\GDHand}{\raisebox{0.1ex}{\scriptsize\color{GDDeepBlue}☞}}

% Override \GDSetHeader to look closer to the screenshot.
\renewcommand{\GDSetHeader}[2]{%
  \pagestyle{fancy}
  \fancyhf{}

  % Top bar with left slashes and right rounded corner
  \fancyhead[C]{%
    \begin{tikzpicture}[baseline=(n.base)]
      \node[inner sep=0pt, outer sep=0pt] (n) {\phantom{X}};
      % main line
      \draw[GDDeepBlue, line width=\GDHeaderRuleWidth, rounded corners=10pt]
        (0,0.22) -- (0.78\linewidth,0.22)
        .. controls +(0.9,0) and +(-0.9,0) .. (0.96\linewidth,0.22)
        -- (\linewidth,0.22);
      % left decorative slashes
      \fill[GDDeepBlue] (0.06\linewidth,0.36) -- (0.08\linewidth,0.36) -- (0.065\linewidth,0.08) -- (0.045\linewidth,0.08) -- cycle;
      \fill[GDDeepBlue!55] (0.09\linewidth,0.36) -- (0.11\linewidth,0.36) -- (0.095\linewidth,0.08) -- (0.075\linewidth,0.08) -- cycle;
    \end{tikzpicture}%
  }

  \fancyhead[L]{\scriptsize\color{GDDeepBlue}#1}
  \fancyhead[R]{\scriptsize\color{GDDeepBlue}#2\;\textbf{\thepage}\;\GDHand}

  % Bottom line + left block + right block (like the sample)
  \fancyfoot[C]{%
    \begin{tikzpicture}[baseline=(n.base)]
      \node[inner sep=0pt, outer sep=0pt] (n) {\phantom{X}};
      \draw[GDDeepBlue, line width=\GDHeaderRuleWidth] (0,0.22) -- (\linewidth,0.22);
      \fill[GDDeepBlue] (0,0.05) rectangle (0.08\linewidth,0.39);
      \fill[GDDeepBlue] (0.88\linewidth,0.05) rectangle (\linewidth,0.39);
    \end{tikzpicture}%
  }
}

% -------------------------
% 13.3) Small answer checkbox at the right side of a box
% Uses tcolorbox overlay (no remember picture) so it never bleeds over other boxes.
% -------------------------
\newcommand{\GDAnswerBoxMark}{%
  \draw[GDBlue, line width=\GDDefaultBoxRule]
    ([xshift=-\GDAnswerMarkOffset, yshift=-\GDAnswerMarkOffset]frame.north east)
    rectangle ++(\GDAnswerMarkSize,-\GDAnswerMarkSize);%
}

% -------------------------
% 13.4) Auto numbering counters
% -------------------------
\newcounter{gdexample}
\newcounter{gdpractice}
\newcounter{gdquestion}
\newcounter{gdproblem}
\newcounter{gddang}

% -------------------------
% 13.5) Easy-to-use boxes (auto-numbered)
% -------------------------
% These are drop-in alternatives so you don't need to type numbers.

% Example (auto): \begin{gdExampleAuto} ... \end{gdExampleAuto}
% Optional arg = number of dotted lines at bottom (default 0 = none)
% e.g. \begin{gdExampleAuto}[5] ... \end{gdExampleAuto}
\newtcolorbox{gdExampleAuto}[1][0]{
  colframe=GDBlue!90!black,
  colback=GDLiteBlue,
  boxrule=\GDExampleBoxRule,
  before upper={\stepcounter{gdexample}},
  title={\GDTag{GDBlue!90!black}{\,\small$\diamond$\;}\;\textbf{Ví dụ\;\arabic{gdexample}}},
  overlay unbroken and last={\GDAnswerBoxMark},
  after upper={\ifnum#1>0\par\vspace{4pt}\DottedLines{#1}\fi},
}

% Practice (auto): \begin{gdPracticeAuto} ... \end{gdPracticeAuto}
\newtcolorbox{gdPracticeAuto}{
  colframe=GDOrange!85!black,
  colback=GDOrange!4,
  boxrule=\GDExampleBoxRule,
  before upper={\stepcounter{gdpractice}},
  title={\GDTag{GDOrange!85!black}{\,\small$\diamond$\;}\;\textbf{Bài\;\arabic{gdpractice}}},
  overlay unbroken and last={\GDAnswerBoxMark},
}

% Question (auto): \begin{gdQuestionAuto} ... \end{gdQuestionAuto}
\newtcolorbox{gdQuestionAuto}{
  colframe=GDPink!85!black,
  colback=GDPink!3,
  boxrule=\GDExampleBoxRule,
  before upper={\stepcounter{gdquestion}},
  title={\GDTag{GDPink!85!black}{\,\small$\diamond$\;}\;\textbf{Câu\;\arabic{gdquestion}}},
  overlay unbroken and last={\GDAnswerBoxMark},
}

% Problem bar (auto): \begin{gdProblemBarAuto} ... \end{gdProblemBarAuto}
\newtcolorbox{gdProblemBarAuto}{
  enhanced,
  breakable,
  colback=white,
  colframe=GDOrange!80!black,
  boxrule=\GDExampleBoxRule,
  left=\GDBoxPadH,right=\GDBoxPadH,top=\GDProblemBarPadV,bottom=\GDProblemBarPadV,
  arc=\GDBoxArc,
  before upper={\stepcounter{gdproblem}},
  title={\GDTag{GDOrange!85!black}{\,\small$\diamond$\;}\;\textbf{Bài\;\arabic{gdproblem}}},
  overlay unbroken and last={\GDAnswerBoxMark},
}

% "Dạng" (auto): \begin{gdDangAuto}{<title>} ... \end{gdDangAuto}
\newtcolorbox{gdDangAuto}[1]{
  colframe=GDPurple,
  colback=white,
  boxrule=\GDDangBoxRule,
  before upper={\stepcounter{gddang}},
  title={\GDTag{GDPurple}{DẠNG\;\arabic{gddang}}\;\textbf{#1}},
}

% -------------------------
% 13.6) Better "Lời giải" block: optional solution + lines
% -------------------------
% Use:
%   \begin{gdSolution}[6]
%     ...optional text...
%   \end{gdSolution}
% If you want ONLY dotted lines, leave the body empty.
\renewenvironment{gdSolution}[1][4]{%
  \def\gdSolutionLines{#1}%
  \LoigiaiTitle
  \vspace{-\GDSolutionVAdjust}
}{%
  \vspace{-\GDSolutionVAdjust}
  \ifnum\gdSolutionLines>0\DottedLines{\gdSolutionLines}\fi
}

